\section{Neighboring Methods and Assumptions}\label{sec:related}

Section~\ref{sec:mergeable-sketches} covers the algebraic lineage:
database algebraic aggregates, MUD aggregation, mergeable summaries,
ordered list homomorphisms, and task-relative sufficient statistics.
Section~\ref{sec:social-measurement} covers the social-science
measurement lineage behind the Manifesto experiment. This section
positions C-TreePO against neighboring methods that share pieces of the
workflow but do not provide the same local-to-global certificate.

\paragraph{Long-Context Methods.}
RAG retrieves relevant chunks but provides no guarantee about what is
\emph{lost} in retrieval; for non-separable targets, retrieval may miss
interaction-bearing information that only becomes decision-relevant when
chunks are combined (Section~\ref{sec:markov-interlude}). Sliding window
methods, recursive summarization, and extended-context models similarly
lack auditable certificates. C-TreePO processes all chunks through a tree
and uses C3-style audits to certify whether cross-chunk interactions are
preserved in the realized pipeline, providing a certification layer that
complements these practical approaches.

\paragraph{Tree-Based Text Composition and Compression.}
Tree-structured inference has been used to compose and compress
natural-language descriptions, with classical work on sentence
compression and fusion providing early structured formulations
\citep{KnightMarcu2000SentenceCompression,BarzilayMcKeown2005SentenceFusion,ClarkeLapata2008ILPSentenceCompression}.
\citet{KuznetsovaEtAl2014TreeTalk} harvest syntactic tree fragments from
web captions and use ILP- and CKY-style structured inference to compose
new image descriptions and generalize captions via parse-tree pruning.
The closest conceptual resonance is that this line already treats
composition and compression together in one tree-based workflow. C-TreePO
differs by making the compressed node a task-relative state and by adding
an auditable local-to-global certificate for oracle preservation and
preference-objective equivalence.

\paragraph{Neural Sufficient Statistics and Simulation-Based Inference.}
The learned-state view of $g$ is closest to neural sufficient-statistic
learning for implicit models. \citet{ChenEtAl2021NASS} learn
approximate sufficient statistics with infomax objectives; the later
sliced sufficient-statistic method of \citet{ChenGutmannWeller2023SSS}
reduces the high-dimensional information problem to many
low-dimensional slice problems. Recent simulation-based inference work
pushes the same idea into downstream likelihood modeling:
\citet{DirmeierAlbertPerezCruz2025SSNL} learn a lower-dimensional
surjective state on which likelihood inference is performed, with
software support in \texttt{sbijax} and \texttt{surjectors}
\citep{DirmeierUlzegaMiraAlbert2024SBIJAX,Dirmeier2024Surjectors}.
Our contribution is the compositional analogue: $g(x)$ is sufficient
only when it preserves all downstream contextual responses
$\fstar(L \concat x \concat R)$. This is why our Markov sketch is a
validation witness rather than an architectural template. Hybrid summary
statistics \citep{MakinenEtAl2024HybridSummaryStatistics} motivate the
same diagnostic discipline: known sketches can be probes or auxiliary
features without defining the general learned state.

\paragraph{Preference Learning and RLHF.}
Direct Preference Optimization fine-tunes a language model from pairwise
human preferences by reparameterizing the Bradley--Terry likelihood so
that the policy itself serves as an implicit reward model
\citep{BradleyTerry1952,RafailovEtAl2023DPO}. Our results apply to DPO
and GRPO variants when the preference signal depends on the input only
through the preserved oracle value. When readability, fluency, or style
preferences fall outside $\fstar$, they must be modeled by a richer
oracle or handled by separate alignment channels.

\paragraph{Assumption-Centered Guarantees.}
Recent machine-learning theory often separates structural invariance
claims from behavioral assumptions on annotator noise and loss geometry.
Our analysis follows that pattern but adds an explicit representational
layer: neural-operator approximation and transfer moduli propose
local-law budgets, C1--C3 and context compatibility drive the
preservation/equivalence results, and quantitative transport to
preference gaps uses explicit method-level Lipschitz and integrability
conditions (Section~\ref{sec:verification}). Keeping this split explicit
makes guarantees easier to audit and stress-test in applied settings.
